\chapter{Descente fid\`element plate}
\label{VIII}

\section{Descente des Modules quasi-coh\'erents}
\label{VIII.1}

Soit $\Sch$
\label{indnot:hb}la cat\'egorie des pr\'esch\'emas. Proc\'edant comme
dans~VI~\ref{VI.11}.b, on trouve que la cat\'egorie des
couples~$(X,F)$ d'un pr\'esch\'ema~$X$ et d'un Module~$F$ sur~$X$,
(o\`u les morphismes sont d\'efinis comme dans \loccit \`a
l'aide de la notion d'image directe de Module par un morphisme
d'espaces annel\'es) peut \^etre consid\'er\'ee comme une
cat\'egorie fibr\'ee au-dessus de~$\Sch$, le foncteur de
changement de base relativement \`a un morphisme~$f\colon X\to Y$
dans $\Sch$ \'etant le foncteur image inverse de Modules
par~$f$. (On notera que la cat\'egorie fibre en~$X\in \Ob(\Sch)$ de
la cat\'egorie fibr\'ee pr\'ec\'edente est la cat\'egorie
\emph{opposée} \`a la cat\'egorie des Modules sur~$X$). Comme
l'image inverse d'un Module quasi-coh\'erent est quasi-coh\'erent,
on voit que la sous-cat\'egorie pleine de la cat\'egorie des
couples~$(X,F)$, form\'ee des couples pour lesquels $F$ est
quasi-coh\'erent, est une sous-cat\'egorie fibr\'ee de la
cat\'egorie fibr\'ee pr\'ec\'edente. (Par contre, si on ne
fait pas d'hypoth\`eses sur~$f$, l'image directe d'un Module
quasi-coh\'erent n'est pas en g\'en\'eral un Module
quasi-coh\'erent). On appellera simplement cette cat\'egorie
fibr\'ee la \emph{catégorie fibrée des Modules
quasi-cohérents sur les préschémas}.

Rappelons d'autre part qu'un morphisme~$f\colon X\to Y$ d'espaces
annel\'es est dit \emph{fidèlement plat}
\index{fidelement plat (morphisme)@fid\`element plat (morphisme)|hyperpage}\index{plat (morphisme fid\`element)|hyperpage}s'il est \emph{plat} (\ie pour tout~$x\in X$, $\mathcal{O}_{X,x}$ est
un module plat sur~$\mathcal{O}_{Y,f(x)}$, \eqref{IV}), et
\emph{surjectif}. On dit que $f$ est un morphisme \emph{quasi-compact}
si l'image inverse par $f$ de toute partie quasi-compacte est
quasi-compacte; lorsque $f$ est un morphisme de pr\'esch\'emas,
cela signifie aussi que l'image inverse par~$f$ d'un ouvert affine
de~$Y$ est r\'eunion \emph{finie} d'ouverts affines de~$X$.

\begin{theoreme}
\label{VIII.1.1}
Soit
$\cal{F}$ la cat\'egorie fibr\'ee des Modules
quasi-coh\'erents sur les pr\'esch\'emas. Soit
$g\colon S'\to S$
un morphisme de pr\'esch\'emas, fid\`element plat et
quasi-compact. Alors $g$ est un morphisme de~$\cal{F}$-descente
effective.
\index{descente effective (morphisme de)|hyperpage}\end{theoreme}

Rappelons\footnote{Nous admettrons ici la th\'eorie
g\'en\'erale de la descente expos\'ee expos\'ee en d\'etail
dans l'article de J.\, \textsc{Giraud} cit\'e dans la
note~\eqref{footnotegiraud} en bas de page de l'Avertissement, travail
que nous citerons \cite{VIII.D} par la suite. \Cf aussi \cite{VIII.2} pour un expos\'e
succinct.} que cela signifie deux choses:

\begin{corollaire}[Descente d'homomorphismes de Modules]
\label{VIII.1.2}
\index{descente d'homomorphismes de modules|hyperpage}Soient $g\colon S'\to S$ un morphisme de pr\'esch\'emas,
fid\`element plat et quasi-compact, $F$ et $G$ deux Modules
quasi-coh\'erents sur~$S$, $F'$ et $G'$ leurs images inverses
sur~$S'$, enfin $F''$ et $G''$ leurs images inverses
sur~$S''=S'\times_{S}S'$. Consid\'erons le diagramme d'applications
d'ensembles d\'efini par les foncteurs changement de base par~$g$,
$p_{1}$, $p_{2}$ (o\`u $\xymatrix@C=.5cm{{p_{1},p_{2}\colon
S'\times_{S}S'}\ar@<2pt>[r]\ar@<-2pt>[r]&{S'}}$ sont les deux
projections):
$$
\xymatrix@C=.5cm{{\mathrm{Hom}_{S}(F,G)}\ar[r]&{\mathrm{Hom}_{S'}(F',G')}\ar@<2pt>[r]\ar@<-2pt>[r]&{\mathrm{Hom}_{S''}(F'',G'').}}
$$
Ce diagramme est exact, \ie d\'efinit une bijection du premier
ensemble sur l'ensemble des co\"incidences des deux applications
\'ecrites du deuxi\`eme dans le troisi\`eme.
\end{corollaire}

En d'autres termes, le foncteur changement de base
par~$g$,
$F\mto F'$, d\'efinit un foncteur \emph{pleinement
fidèle} de la cat\'egorie des Modules quasi-coh\'erents sur~$S$ dans la cat\'egorie des Modules quasi-coh\'erents sur~$S'$
munis d'une donn\'ee de descente relativement
\`a~$g$.
De plus:

\begin{corollaire}[Descente de Modules]
\label{VIII.1.3}
\index{descente de modules|hyperpage}Pour tout Module quasi-coh\'erent~$F'$ sur~$S'$, toute donn\'ee de
descente sur~$F'$ relativement \`a $g$ est \emph{effective},
\index{donn\'ee de descente effective|hyperpage}\ie $F'$ est isomorphe avec sa donn\'ee de descente \`a l'image
inverse par~$g$ d'un Module quasi-coh\'erent sur~$S$
(d\'etermin\'e \`a isomorphisme unique pr\`es en vertu
de~\ref{VIII.1.2}).
\end{corollaire}

En d'autres termes, le foncteur pleinement fid\`ele
pr\'ec\'edent est m\^eme une
\emph{équivalence}. Pratiquement, cela signifie qu'il revient au
m\^eme de se donner un Module quasi-coh\'erent sur~$S$, ou un
Module quasi-coh\'erent sur~$S'$ muni d'une donn\'ee de descente
relativement \`a~$g$.

\subsubsection*{D\'emonstration de~\ref{VIII.1.1}}
Soit
d'abord $T$ un~$S$-pr\'esch\'ema qui est $S$-isomorphe \`a
la somme d'une famille
d'ouverts
induits~$S_{i}$ de~$S$ qui recouvrent~$S$. Alors il est \'evident
que le morphisme structural~$T\to S$ est un morphisme de
$\cal{F}$-descente effective (cela signifie pr\'ecis\'ement que la
donn\'ee d'un Module quasi-coh\'erent~$F$ sur~$S$ \'equivaut
\`a la donn\'ee de Modules quasi-coh\'erents~$F_{i}$ sur
les~$S_{i}$, et d'isomorphismes de recollement~$\varphi_{ji}\colon
F_{i}|S_{i}\cap S_{j}\to F_{j}|S_{i}\cap S_{j}$ satisfaisant la
condition de cocha\^ines bien connue). En vertu de~VII,~8
il s'ensuit que pour v\'erifier que $g\colon S'\to S$ est un
morphisme de~$\cal{F}$-descente effective, il suffit de le
v\'erifier pour le morphisme~$g_{T}\colon T'=T\times_{S}S'\to T$
d\'eduit de~$g$ par le changement de base~$T\to S$. (Remarquer que
l'hypoth\`ese sur~$T\to S$ reste stable par changement de base
quelconque, donc que $T\to S$ est en fait un morphisme de
$\cal{F}$-descente effective \emph{universel}). Prenant pour~$S_{i}$
des ouverts affines qui recouvrent~$S$, on est donc ramen\'e au cas
o\`u~$S$ est affine.

Alors $S'$ est r\'eunion finie d'ouverts affines, et prenant le
$S$-sch\'ema somme de ces derniers, on trouve un $S$-sch\'ema
affine~$S_{1}$ et un $S$-morphisme~$S_{1}\to S'$ plat et
surjectif. Donc $S_{1}$ est aussi fid\`element plat sur~$S$. Si donc
on prouve qu'un morphisme fid\`element plat et affine est un
morphisme de~$\cal{F}$-descente effective, donc un morphisme de
$\cal{F}$-descente strict universel, (l'hypoth\`ese \'etant en
effet stable par changement de base), on en conclut en particulier que
le morphisme structural~$S_{1}\to S$ est un morphisme de
$\cal{F}$-descente strict universel, et comme il existe un
$S$-morphisme~$S_{1}\to S'$, il en r\'esultera bien, par \cite{VIII.D}, que
$g\colon S'\to S$ est un morphisme de~$\cal{F}$-descente strict.

\enlargethispage{.5cm}Cela nous ram\`ene donc au cas o\`u $g$ est un morphisme affine,
et comme on a vu on peut alors de plus supposer $S$ affine, donc
\emph{on peut supposer $S$ et $S'$ affines}. Dans ce cas,
\ref{VIII.1.2} \'equivaut au

\begin{lemme}
\label{VIII.1.4}
Soient $A$ un anneau, $A'$ une~$A$-alg\`ebre fid\`element plate,
$M$ et $N$ deux~$A$-modules, $M'$ et $N'$ les~$A'$-modules d\'eduits
par changement d'anneau~$A\to A'$, et $M''$,$N''$
les~$A''=A'\otimes_{A}A'$-modules d\'eduits par changement
d'anneau~$A\to A''$. Alors la suite d'applications ensemblistes
$$
\xymatrix@C=.5cm{{\mathrm{Hom}_{A}(M,N)}\ar[r]&{\mathrm{Hom}_{A'}(M',N')}\ar@<2pt>[r]\ar@<-2pt>[r]&{\mathrm{Hom}_{A''}(M'',N'')}}
$$
est exacte.
\end{lemme}

Comme
l'homomorphisme~$N\to N'$ est injectif ($A'$ \'etant fid\`element
plat sur~$A$) on voit que la premi\`ere fl\`eche est injective. Il
reste \`a prouver que si un
$A'$-homomorphisme
$u'\colon M'\to N'$ est compatible avec les donn\'ees de descente,
alors il provient d'un $A$-homomorphisme~$u\colon M\to N$. Or cela
signifie aussi simplement que $u'$ applique le sous-ensemble~$M$
de~$M'$ dans le sous-ensemble~$N$ de~$N'$ (l'application~$u\colon M\to
N$ induite sera alors automatiquement $A$-lin\'eaire puisque $u'$
est $A'$-lin\'eaire, et on voit de m\^eme que $u'$ est
n\'ecessairement \'egal \`a~$u\otimes_{A}A'$). Or si $x\in M$,
alors $u'(x)$ est un \'el\'ement dans le noyau du couple
d'applications~$\xymatrix@C=.5cm{{N'}\ar@<2pt>[r]\ar@<-2pt>[r]&{N''}}$. On est
donc ramen\'e pour prouver~\ref{VIII.1.4}
au cas particulier suivant (correspondant au cas o\`u~$M=A$):

\begin{corollaire}
\label{VIII.1.5}
Soit $N$ un~$A$-module, alors la suite d'applications ensemblistes
$$
\xymatrix@C=.5cm{{N}\ar[r]&{N'}\ar@<2pt>[r]\ar@<-2pt>[r]&{N''}}
$$
est exacte.
\end{corollaire}

Soit en effet $A_{1}$ une~$A$-alg\`ebre fid\`element plate. Pour
montrer que la suite envisag\'ee est exacte, il suffit de prouver
que la suite qui s'en d\'eduit par le changement d'anneau~$A\to
A_{1}$ l'est. Or cette derni\`ere, comme on voit de suite, est celle
relative au
$A_1$-module
$N_{1}=N\otimes_{A}A_{1}$ et \`a la $A_{1}$-alg\`ebre
$A_{1}'=A_{1}\otimes_{A}A'$. Il suffit donc de trouver un $A_{1}$
fid\`element plat sur~$A$, tel que $\Spec(A_{1}')\to \Spec(A_{1})$
soit un morphisme de~$\cal{F}$-descente strict. Or il suffit en effet
de prendre $A_{1}=A'$, car alors le morphisme pr\'ec\'edent admet
un morphisme inverse \`a droite, donc en vertu de~\cite{VIII.D} c'est un
morphisme de descente effective pour n'importe quelle cat\'egorie
fibr\'ee sur~$\Sch$.

Il reste enfin \`a montrer que si $N'$ est un~$A'$-module muni d'une
donn\'ee de descente pour~$A\to A'$, \ie muni d'un isomorphisme
$$
\varphi\colon N_{1}'\isomto N_{2}'
$$
entre les deux modules d\'eduits de~$N$ par les changements
d'anneaux $\xymatrix@C=.5cm{{A'}\ar@<2pt>[r]\ar@<-2pt>[r]&{A'\otimes_{A}A'}}$,
alors
$N'$ est isomorphe avec sa donn\'ee de descente \`a un
module~$N\otimes_{A}A'$. On voit facilement, compte tenu
de~\ref{VIII.1.5}, que cet \'enonc\'e \'equivaut au suivant:

\begin{lemme}
\label{VIII.1.6}
Soit $N'$ un~$A'$-module muni d'une donn\'ee de descente
relativement \`a~$A\to A'$ (o\`u $A'$ est une
$A$-alg\`ebre). Soit $N$ le sous-$A$-module de~$N'$ form\'e
des~$x$ tels que $\varphi(x\otimes_{A'}1_{A'})=1_{A'}\otimes_{A'}x$,
et consid\'erons l'homomorphisme canonique
$$
N\otimes_{A}A'\to N'\quoi,
$$
(qui est alors compatible avec les donn\'ees de descente). Si $A'$
est fid\`element plat sur~$A$, cet homomorphisme est un
isomorphisme.
\end{lemme}

D\'emontrons ce lemme. Soit encore $A_{1}$ une $A$-alg\`ebre
fid\`element plate, pour montrer que le morphisme envisag\'e est
un isomorphisme, il suffit de prouver qu'il le devient apr\`es le
changement d'anneau
$A\to A_1$.
Or, utilisant la platitude de~$A_{1}$ sur~$A$, on voit que
l'homomorphisme ainsi obtenu n'est autre que celui qu'on obtiendrait
directement en termes du module~$N'\otimes_{A}A_{1}$ sur~$A_{1}'=A'\otimes_{A}A_{1}$, muni de la donn\'ee de descente
relativement \`a~$A_{1}\to A_{1}'$ qui se d\'eduit canoniquement
par changement d'anneau de celle qui \'etait donn\'ee
sur~$N'$. Ainsi il suffit de trouver un~$A_{1}$ fid\`element plat
sur~$A$ tel que $\Spec(A_{1}')\to \Spec(A_{1})$ soit un morphisme de
$\cal{F}$-descente effective. On prend alors comme ci-dessus
$A_{1}=A'$. Cela ach\`eve la d\'emonstration de~\ref{VIII.1.6}, et
par l\`a la d\'emonstration de~\ref{VIII.1.1}.

\begin{corollaire}[Descente de sections de Modules]
\label{VIII.1.7}
\index{descente de sections de modules|hyperpage}Soit $g\colon S'\to S$ un morphisme de pr\'esch\'emas,
fid\`element plat et quasi-compact. Pour tout Module
quasi-coh\'erent $G$ sur~$S$, soient $G'$ et $G''$ ses images
inverses sur~$S'$ et $S''=S'\times_{S}S'$, et consid\'erons le
diagramme d'homomorphismes de Modules sur~$S$:
$$
\xymatrix@C=.5cm{{G}\ar[r]&{g_*(G')}\ar@<2pt>[r]\ar@<-2pt>[r]&{h_*(G'')}}
$$
(o\`u $h\colon S''\to S$ est le morphisme structural). Ce diagramme
est \emph{exact}.
\end{corollaire}

En
effet, cela signifie que pour tout ouvert~$U$ dans~$S$, le diagramme
correspondant form\'e par les sections sur~$U$ est exact. On peut
\'evidemment supposer alors $U=S$, et l'exactitude en question est
alors un cas particulier de~\ref{VIII.1.2}, obtenu en faisant
$F=\mathcal{O}_{S}$.

Comme le foncteur image inverse de Modules est exact \`a droite, on
conclut formellement \`a partir de~\ref{VIII.1.1}:

\begin{corollaire}[Descente de Modules quotients]
\label{VIII.1.8}
\index{descente de modules quotients|hyperpage}Avec les notations de~\ref{VIII.1.7}, soit de plus, pour tout Module
quasi-coh\'erent~$F$ sur un pr\'esch\'ema, $\Quot(F)$ l'ensemble
des Modules quasi-coh\'erents quotients de~$F$. Avec cette
convention, le diagramme d'applications d'ensembles:
$$
\xymatrix@C=.5cm{{\mathrm{Quot}(G)}\ar[r]&{\mathrm{Quot}(G')}\ar@<2pt>[r]\ar@<-2pt>[r]&{\mathrm{Quot}(G'')}}
$$
est exact.
\end{corollaire}

(On aurait \'evidemment le m\^eme \'enonc\'e avec les
sous-Modules au lieu de Modules quotients, puisque les deux se
correspondent biunivoquement). Faisant en particulier
$G=\mathcal{O}_{S}$, on trouve:

\begin{corollaire}[Descente des sous-pr\'esch\'emas ferm\'es]
\label{VIII.1.9}
\index{descente des sous-pr\'esch\'emas ferm\'es|hyperpage}Pour tout pr\'esch\'ema~$X$, soit $H(X)$ l'ensemble des
sous-pr\'esch\'emas ferm\'es de~$X$. Avec cette notation, et
sous les conditions de~\ref{VIII.1.7}, le diagramme d'applications
d'ensembles suivant
$$
\xymatrix@C=.5cm{{H(S)}\ar[r]&{H(S')}\ar@<2pt>[r]\ar@<-2pt>[r]&{H(S'')}}
$$
\emph{est exact}.
\end{corollaire}

Il y a lieu de compl\'eter le th\'eor\`eme~\ref{VIII.1.1} par le
r\'esultat suivant:

\begin{proposition}[Descente de propri\'et\'es de Modules]
\label{VIII.1.10}
\index{descente de propri\'et\'es de modules|hyperpage}Soient $g\colon S'\to S$ un morphisme fid\`element plat et
quasi-compact, $F$ un Module quasi-coh\'erent sur~$S$. Pour que~$F$
soit de type fini, \resp de pr\'esentation finie, \resp localement
libre et de type fini, il faut et il suffit que son image inverse $F'$
sur~$S'$ le soit.
\end{proposition}

Il n'y a qu'\`a prouver le \og il suffit\fg. On peut \'evidemment
supposer~$S$ affine,
et rempla\c cant alors~$S'$ par
une
somme d'ouverts affines recouvrant~$S'$, on est ramen\'e au cas
o\`u~$S'$ est \'egalement affine. Alors notre \'enonc\'e
\'equivaut au suivant:

\begin{corollaire}
\label{VIII.1.11}
Soient $A$ un anneau, $A'$ une~$A$-alg\`ebre fid\`element plate,
$M$ un~$A$-module, $M'$ le~$A'$-module $M\otimes_{A}A'$. Pour que~$M$
soit de type fini (\resp de pr\'esentation finie, \resp localement
libre de type fini) il faut et il suffit que~$M'$ le soit.
\end{corollaire}

En effet, on a $M=\varinjlim_i M_{i}$, o\`u les~$M_{i}$
sont les sous-modules de type fini de~$M$. Par suite
$M'=\varinjlim_i M_{i}'$, et si $M'$ est de type fini, $M'$
est \'egal \`a l'un des~$M_{i}'$, donc par fid\`ele platitude
$M$ est \'egal \`a~$M_{i}$, donc $M$ est de type fini. Par suite
il existe une suite exacte
$$
0\to R\to L\to M\to 0\quoi,
$$
avec $L$ libre de type fini, d'o\`u une suite exacte
$$
0\to R'\to L'\to M'\to 0\quoi,
$$
avec $L'$ libre de type fini. Si donc $M'$ est de pr\'esentation
finie, $R'$ est de type fini, donc d'apr\`es ce qui pr\'ec\`ede
$R$ est de type fini, donc $M$ est de pr\'esentation finie. Enfin,
dire que~$M$ est localement libre et de type fini, signifie qu'il est
de pr\'esentation finie et plat (\cf \ref{IV} dans le cas
noeth\'erien; le cas g\'en\'eral est laiss\'e au
lecteur). Comme chacune de ces propri\'et\'es se descend bien, il
en est de m\^eme de leur conjonction, ce qui ach\`eve la
d\'emonstration.

\begin{remarque}
\label{VIII.1.12}
La conjonction de~\ref{VIII.1.1} et de~\ref{VIII.1.10} montre que
l'\'enonc\'e~\ref{VIII.1.1} reste encore valable, quand on
remplace la cat\'egorie fibr\'ee~$\cal{F}$ par la
sous-cat\'egorie fibr\'ee form\'ee des Modules
quasi-coh\'erents de type fini, \resp de pr\'esentation finie,
\resp localement libres de type fini, \resp localement libres de rang
donn\'e $n$.
\end{remarque}

\section{Descente des pr\'esch\'emas affines sur un autre}
\label{VIII.2}

Comme le foncteur image inverse de Modules est compatible avec le
produit tensoriel et d'autres op\'erations tensorielles, le
th\'eor\`eme~\ref{VIII.1.1} implique diverses variantes, obtenues
en envisageant, au lieu d'un seul Module quasi-coh\'erent, un Module
quasi-coh\'erent ou un syst\`eme de Modules quasi-coh\'erents
muni de structures suppl\'ementaires diverses s'exprimant \`a
l'aide d'op\'erations tensorielles. Par exemple, la donn\'ee de
trois Modules quasi-coh\'erents~$F$, $G$, $H$ sur~$S$ et d'un
accouplement
$$
F\otimes G\to H\quoi,
$$
\'equivaut \`a la donn\'ee de trois Modules
quasi-coh\'erents~$F'$, $G'$, $H'$ sur~$S'$, munis de donn\'ees de
descente relativement \`a~$g\colon S'\to S$, et munis d'un
accouplement
$$
F'\otimes G'\to H'\quoi,
$$
\og compatible\fg avec ces donn\'ees de descente, au sens \'evident
du terme. Par exemple, si $F=G=H$, on voit que la donn\'ee d'un
Module quasi-coh\'erent~$F$ sur~$S$ muni d'une loi d'alg\`ebre
(dont pour l'instant nous ne supposons pas qu'elle satisfasse \`a
aucun axiome d'associativit\'e, commutativit\'e ou d'existence
d'une section unit\'e), \'equivaut \`a la m\^eme donn\'ee
sur~$S'$, munie en plus d'une donn\'ee de descente. Et utilisant les
r\'esultats du num\'ero pr\'ec\'edent, on constate
aussit\^ot que pour que~$F$ satisfasse \`a l'un des axiomes
habituels auxquels on vient de faire allusion, il faut et il suffit
qu'il en soit ainsi pour~$F'$. Par exemple, la donn\'ee d'une
Alg\`ebre quasi-coh\'erente~$\mathcal{A}$ sur~$S$ (par quoi nous
sous-entendons d\'esormais: associative, commutative, \`a section
unit\'e) \'equivaut \`a la donn\'ee d'une Alg\`ebre
quasi-coh\'erente~$\mathcal{A}'$ sur~$S'$, munie d'une donn\'ee de
descente relativement \`a
$g\colon S'\to S$.
Si on se rappelle l'\'equivalence entre la cat\'egorie duale des
Alg\`ebres quasi-coh\'erentes sur~$S$, et de la cat\'egorie
des~$S$-pr\'esch\'emas affines sur~$S$, (EGA~II par.1), on trouve
aussit\^ot:

\begin{theoreme}
\label{VIII.2.1}
Soit $\cal{F}'$ la cat\'egorie fibr\'ee des morphismes affines de
pr\'esch\'emas $f\colon X\to S$, consid\'er\'ee comme
sous-cat\'egorie fibr\'ee de la cat\'egorie fibr\'ee
des fl\`eches dans la cat\'egorie des pr\'esch\'emas $\Sch$
\textup{(VI~\ref{VI.11}.a)}. Soit $g\colon S'\to S$ un morphisme de
pr\'esch\'emas fid\`element plat et quasi-compact. Alors $g$ est
un morphisme de~$\cal{F}'$-descente effective.
\end{theoreme}

\section[Descente de propri\'et\'es ensemblistes]{Descente de propri\'et\'es ensemblistes et de propri\'et\'es de finitude de morphismes\footnotemark}
\label{VIII.3}


\footnotetext{Pour d'autres r\'esultats comme ceux trait\'es dans les
num\'eros 3 et~4, \Cf EGA~IV~2.3, 2.6,~2.7.}

\begin{proposition}
\label{VIII.3.1}
Soient $f\colon X\to Y$ un~$S$-morphisme, $g\colon S'\to S$ un
morphisme surjectif, $f'\colon X'=X\times_{S}S'\to Y'=Y\times_{S}S'$
le morphisme d\'eduit de~$f$ par le changement de base \`a l'aide
de~$g\colon S'\to S$. Pour que $f$ soit surjectif (\resp radiciel), il
faut et il suffit que $f'$ le soit.
\end{proposition}

On note que $f'$ peut aussi s'obtenir par le changement de base~$Y'\to
Y$, qui est \'egalement surjectif puisque d\'eduit de~$g\colon
S'\to S$ qui l'est. D'autre part, pour tout~$y\in Y$ et tout~$y'\in Y'$
au-dessus de~$y$, on a un isomorphisme
$$
X'_{y'}\isomfrom X_{y}\otimes_{\kres(y)}\kres(y')\quoi,
$$
o\`u $X_{y}$ d\'esigne la fibre de~$X$~en~$y$, et $X'_{y'}$ celle
de~$X'$~en $y'$. Il en r\'esulte que~$X_{y}$ est non vide (\resp a
au plus un point, et ce dernier correspond \`a une extension
r\'esiduelle radicielle) si et seulement si
$X'_{y'}$
a la m\^eme propri\'et\'e. Cela prouve~\ref{VIII.3.1}.

\begin{corollaire}
\label{VIII.3.2}
Sous les conditions de~\ref{VIII.3.1}, si $f'$ est injectif
(\resp bijectif), alors $f$ l'est \'egalement.
\end{corollaire}

Cela provient du fait que si $X'_{y'}$ a au plus un point
(\resp exactement un point) il en est de m\^eme de~$X_{y}$; il en
est bien ainsi, puisque le morphisme $X'_{y'}\to X_{y}$ est surjectif
(car d\'eduit de~$\Spec(\kres(y'))\to \Spec(\kres(y))$ qui l'est).

\begin{proposition}
\label{VIII.3.3}
Avec les notations de~\ref{VIII.3.1}. Supposons que~$g\colon S'\to S$
soit surjectif et quasi-compact (\resp fid\`element plat et
quasi-compact). Pour que~$f$ soit quasi-compact (\resp de type fini),
il faut et il suffit que~$f'$ le soit.
\end{proposition}

Il
y a \`a prouver seulement le~\og il suffit\fg. On peut \'evidemment
supposer~$S=Y$, puisque l'hypoth\`ese faite sur~$g\colon S'\to S$
est conserv\'ee pour~$Y'\to Y$. De plus on peut supposer~$Y$
affine. Alors $Y'$ est quasi-compact, donc $X'$ est quasi-compact
(puisque~$f'$ l'est par hypoth\`ese). Soit $(X_{i})_{i\in I}$ une
famille d'ouverts affines de~$X$ recouvrant~$X$, alors les~$X_{i}'$
sont des ouverts de~$X'$ recouvrant~$X'$, donc il y a une sous-famille
finie qui recouvre~$X'$. Comme $X'\to X$ est surjectif, il s'ensuit
que les~$X_{i}$ correspondants recouvrent d\'ej\`a~$X$, donc~$X$ est
quasi-compact, \ie $f$ est quasi-compact. Supposons maintenant $f'$
de type fini, prouvons que~$f$ l'est, en supposant~$g$ fid\`element
plat. Rempla\c cant $Y'$ par la somme d'une famille d'ouverts
affines le recouvrant, on peut supposer $Y'$ affine. Enfin, $X$
\'etant recouvert par un nombre fini d'ouverts affines $X_{i}$ par
ce qui pr\'ec\`ede, il faut montrer qu'ils sont de type fini sur~$Y$ sachant que $X_{i}'$ est de type fini sur~$Y'$. Cela nous
ram\`ene alors au

\begin{corollaire}
\label{VIII.3.4}
Soient $B$ une $A$-alg\`ebre, $A'$ une~$A$-alg\`ebre
fid\`element plate, $B'=B\otimes_{A}A'$ la~$A$-alg\`ebre
d\'eduite de~$B$ par changement d'anneau. Pour que~$B$ soit de type
fini, il faut et il suffit que~$B'$ le soit.
\end{corollaire}

Il y a \`a prouver seulement le~\og il suffit\fg. On a
$B=\varinjlim_i B_{i}$, o\`u les $B_{i}$ sont les
sous-alg\`ebres de type fini de~$B$. On a donc
$B'=\varinjlim_i B_{i}'$, et si $B'$ est de type fini
sur~$A'$, alors $B'$ est \'egal \`a un des~$B_{i}'$, donc $B$ est
\'egal \`a~$B_{i}$, donc est de type fini.

\begin{corollaire}
\label{VIII.3.5}
Supposons encore le morphisme de changement de base~$g\colon S'\to S$
fid\`element plat et quasi-compact. Pour que~$f$ soit quasi-fini, il
faut et il suffit que~$f'$ le soit.
\end{corollaire}

En effet, la propri\'et\'e \og quasi-fini\fg est par d\'efinition
la conjonction de \og type fini\fg et \og \`a fibres finies\fg, dont
chacune se descend bien par~$g$, la premi\`ere en vertu
de~\ref{VIII.3.3}, la seconde par le raisonnement de~\ref{VIII.3.1}
(n'utilisant que le fait que $g$ soit surjectif).

\begin{remarques}
\label{VIII.3.6}
Soient $A$ un anneau, $X$ un~$A$-pr\'esch\'ema. On voit facilement
que
les conditions suivantes sont \'equivalentes:

\begin{enumerate}
\item[(i)] Il existe un anneau noeth\'erien~$A_{0}$ (qu'on peut si
on veut supposer un sous-anneau de type fini de~$A$), un
$A_{0}$-pr\'esch\'ema de type fini~$X_{0}$, un homomorphisme
$A_0\to A$
et un $A$-isomorphisme~$X\isomto X_{0}\times_{A_{0}}A$.
\item[(ii)] Le morphisme diagonal~$X\to X\times_{\Spec(A)}X$ est
quasi-compact (condition vide si $X$ est s\'epar\'e sur~$A$), $X$
est r\'eunion finie d'ouverts affines~$X_{i}$ dont les
anneaux~$B_{i}$ sont des alg\`ebres de pr\'esentation finie
sur~$A$, \ie quotients d'alg\`ebres de polyn\^omes \`a un
nombre fini d'ind\'etermin\'ees, par des id\'eaux de type fini.
\end{enumerate}

Si $X$ est lui-m\^eme affine, d'anneau~$B$, ces conditions
signifient simplement que~$B$ est une alg\`ebre de pr\'esentation
finie sur~$A$.

Un morphisme~$f\colon X\to Y$ est dit \emph{morphisme de
présentation finie},
\index{presentation finie@pr\'esentation finie (morphisme de)|hyperpage}et on dit encore que~$X$ est de pr\'esentation finie sur~$Y$, si $Y$
est r\'eunion d'ouverts affines~$Y_{i}$, tel que~$X|Y_{i}$ en tant
que $Y_{i}$-pr\'esch\'ema satisfasse aux conditions
\'equivalentes pr\'ec\'edentes. Il en est alors de m\^eme pour
$X|Y'$ pour \emph{tout} ouvert affine $Y'$ dans $Y$. C'est l\`a une
propri\'et\'e stable par changement de base, et d'ailleurs le
compos\'e de deux morphismes de pr\'esentation finie est de
pr\'esentation finie.

Ces notions pos\'ees, on voit sur (ii), proc\'edant comme
dans~\ref{VIII.1.10}, que cet \'enonc\'e reste valable en y
rempla\c cant les mots \og de type fini\fg par \og de pr\'esentation
finie\fg.
\end{remarques}

\section{Descente de propri\'et\'es topologiques}
\label{VIII.4}

\begin{theoreme}
\label{VIII.4.1}
Soient $g\colon Y'\to Y$ un morphisme, et $Z$ une partie de~$Y$. On
suppose que~$g$ est plat, et qu'il existe un morphisme
quasi-compact~$f\colon X\to Y$ tel que~$Z=f(X)$ (N.B. si $Y$ est
noeth\'erien, cette derni\`ere condition est impliqu\'ee par
\og $Z$ est constructible\fg). Alors on a
$$
g^{-1}(\overline{Z})=\overline{g^{-1}(Z)}\quoi.
$$
\end{theoreme}

On peut supposer~$Y$ affine, puis $Y'$ affine. Comme $Y$ est affine,
$X$ est r\'eunion
finie d'ouverts affines~$X_{i}$, et rempla\c cant~$X$ par la somme
des~$X_{i}$, on peut supposer \'egalement~$X$
affine. Soient~$A,A',B$ les anneaux de~$Y,Y',X$, $B'=B\otimes_{A}A'$
celui de~$X'=X\times_{Y}Y'$, $I$ le noyau de~$A\to B$, $I'$ le noyau
de~$A'\to B'$, donc les parties ferm\'ees de~$Y$ et~$Y'$
d\'efinies par ces id\'eaux sont respectivement l'adh\'erence
de~$Z=f(X)$ et l'adh\'erence de~$Z'=f'(X')=g^{-1}(Z)$. On veut
\'etablir que cette derni\`ere est \'egale \`a
$g^{-1}(\overline{Z})$, ce qui r\'esultera de~$I'=IA'$, lui-m\^eme
cons\'equence de la platitude de~$A'$ sur~$A$.

\begin{corollaire}
\label{VIII.4.2}
Soient $g\colon Y'\to Y$ un morphisme plat et quasi-compact, et $Z'$
une partie ferm\'ee de~$Y'$ satur\'ee pour la relation
d'\'equivalence ensembliste d\'efinie par $g$. Alors on a
$Z'=g^{-1}(\overline{g(Z')})$.
\end{corollaire}

\enlargethispage{.5cm}On a en effet $Z'=g^{-1}(Z)$, avec~$Z=g(Z')$. On peut alors
appliquer~\ref{VIII.4.1}, en notant que la condition mise sur~$Z$
dans~\ref{VIII.4.1} est bien v\'erifi\'ee en prenant pour~$X$ le
pr\'esch\'ema~$Z'$ muni de la structure r\'eduite induite
par~$Y'$. (Le fait que~$g$ soit quasi-compact assure alors que~le
morphisme induit~$f\colon Z'\to Y$ est quasi-compact).

L'\'enonc\'e~\ref{VIII.4.2} signifie aussi que \emph{la topologie
de~$g(Y')$ induite par~$Y$ est quotient de celle de~$Y'$}. En
particulier:

\begin{corollaire}
\label{VIII.4.3}
Soit $g\colon Y'\to Y$ un morphisme fid\`element plat et
quasi-compact. Alors $g$ fait de~$Y$ un espace topologique quotient
de~$Y'$, \ie pour une partie~$Z$ de~$Y$, $Z$ est ferm\'ee
(\resp ouverte) si et seulement si $Z'=g^{-1}(Z)$ l'est.
\end{corollaire}

Rappelons maintenant que deux \'el\'ements $a,b$, de~$Y'$ ont
m\^eme image dans $Y$ si et seulement si ils sont de la forme
$p_{1}(c),p_{2}(c)$ pour un \'el\'ement convenable $c$
dans~$Y''=Y'\times_{Y}Y'$. Il en r\'esulte que si $g$ est surjectif,
on a un diagramme \emph{exact} d'ensembles
$$
\xymatrix@C=.5cm{{\cal{P}(Y)}\ar[r]&{\cal{P}(Y')}\ar@<2pt>[r]\ar@<-2pt>[r]&{\cal{P}(Y'')}}\quoi,
$$
o\`u pour tout ensemble~$E$, on d\'esigne par~$\cal{P}(E)$
l'ensemble de ses parties. Ceci pos\'e, \ref{VIII.4.3} peut aussi
s'interpr\'eter ainsi:

\begin{corollaire}[Descente des parties ouvertes \resp ferm\'ees]
\label{VIII.4.4}
Soit
\index{descente des parties ouvertes (\resp ferm\'ees)|hyperpage}$g\colon Y'{\to} Y$ comme dans~\ref{VIII.4.3}. Pour tout
pr\'esch\'ema~$X$, soit $\Ouv(X)$ \resp $\Fer(X)$ l'ensemble de
ses parties ouvertes \resp l'ensemble de ses parties
ferm\'ees. Alors on a des diagrammes \emph{exacts} d'applications
ensemblistes (d\'eduits de~$g$ et des deux projections de
$Y''=Y'\times_{Y}Y'$):
\vspace*{-5pt}
\begin{gather*}
\xymatrix@C=.5cm{{\mathrm{Ouv}(Y)}\ar[r]&{\mathrm{Ouv}(Y')}\ar@<2pt>[r]\ar@<-2pt>[r]&{\mathrm{Ouv}(Y'')}}\\[-3pt]
\xymatrix@C=.5cm{{\mathrm{Fer}(Y)}\ar[r]&{\mathrm{Fer}(Y')}\ar@<2pt>[r]\ar@<-2pt>[r]&{\mathrm{Fer}(Y'')}}
\end{gather*}
\end{corollaire}

On a le compl\'ement suivant \`a~\ref{VIII.4.3}:

\begin{corollaire}
\label{VIII.4.5}
Soit $g\colon Y'\to Y$ comme dans~\ref{VIII.4.3}, et soit~$Z$ une
partie de~$Y$ telle qu'il existe un morphisme quasi-compact~$f\colon
X\to Y$ d'image~$Z$ (par exemple $Z$ constructible, $Y$
noeth\'erien). Pour que $Z$ soit une partie localement ferm\'ee
de~$Y$, il faut et il suffit que $Z'=g^{-1}(Z)$ soit une partie
localement ferm\'ee de~$Y'$.
\end{corollaire}

Il suffit de prouver le \og il suffit\fg. Soit $Y_{1}$ le
sous-pr\'esch\'ema ferm\'e de~$Y$, adh\'erence de~$Z$ muni de
la structure r\'eduite induite, et soit~$Y_{1}'=Y_{1}\times_{Y}Y'$
le sous-pr\'esch\'ema ferm\'e de~$Y'$ image r\'eciproque
de~$Y_{1}$. Son ensemble sous-jacent est l'image
inverse~$g^{-1}(Y_{1})=g^{-1}(\overline{Z})$, donc est \'egal en
vertu de~\ref{VIII.4.1} \`a~$\overline{Z'}$. Comme $Z'$ est
localement ferm\'e dans~$Y'$, il est ouvert dans~$\overline{Z'}$
donc ouvert dans~$Y_{1}'$. Or ce dernier est fid\`element plat et
quasi-compact sur~$Y_{1}$, donc en vertu de~\ref{VIII.4.3} on en
conclut que~$Z$ est ouvert dans~$Y_{1}$, \ie dans~$\overline{Z}$, ce
qui signifie que $Z$ est localement ferm\'e.

\begin{corollaire}
\label{VIII.4.6}
Soit $g\colon S'\to S$ un morphisme fid\`element plat et
quasi-compact, $f\colon X\to Y$ un $S$-morphisme, $f'\colon X'\to Y'$
le $S'$-morphisme qui s'en d\'eduit par changement de
base. Supposons que~$f'$ soit une application ouverte (\resp une
application ferm\'ee, \resp quasi-compacte
et un hom\'eomorphisme dans, \resp un hom\'eomorphisme sur); alors
$f$ a la m\^eme propri\'et\'e.
\end{corollaire}

Comme $Y'$ est fid\`element plat et quasi-compact sur~$Y$, on peut
supposer~$Y=S$. Soit~$Q$
une partie de~$X$, on a alors (en d\'esignant par $h$ le morphisme
de projection $X'\to X$):
$$
g^{-1}(f(Q))=f'(h^{-1}(Q))\quoi.
$$

Si $Q$ est ouvert (\resp ferm\'e), il en est de m\^eme
de~$h^{-1}(Q)$, donc aussi de~$f'(h^{-1}(Q))$ si on suppose $f'$ une
application ouverte (\resp ferm\'ee), donc il en est de m\^eme
de~$f(Q)$ en vertu de la formule pr\'ec\'edente et
de~\ref{VIII.4.3}. Cela prouve les deux premi\`eres assertions
dans~\ref{VIII.4.6}, il reste \`a examiner le cas o\`u~$f'$ est un
hom\'eomorphisme dans, et prouver alors que~$f$ est un
hom\'eomorphisme dans. (Le cas d'un hom\'eomorphisme sur
r\'esultera alors de~\ref{VIII.3.1}). En vertu
de~\ref{VIII.3.2}
$f$ est injectif, il reste \`a prouver que l'application $X\to f(X)$
est ouverte. On sait d\'ej\`a que~$f$ est quasi-compact en vertu
de~\ref{VIII.3.3}. Il suffit de prouver que pour toute partie
ferm\'ee~$Z$ de~$X$, on a $Z=f^{-1}(\overline{f(Z)})$, ce qui
\'equivaut \`a la formule analogue pour les images inverses par
l'application surjective~$h\colon X'\to X$, \ie \`a
$$
Z'=f'^{-1}(g^{-1}(\overline{f(Z)})\quoi,
$$
o\`u on pose $Z'=h^{-1}(Z)$. Or en vertu de~\ref{VIII.4.1}
appliqu\'e \`a la partie~$f(Z)$ de~$Y$, on a
$g^{-1}(\overline{f(Z)})=\overline{g^{-1}(f(Z))}$, et la formule \`a
prouver \'equivaut \`a
$$
Z'=f'^{-1}(\overline{f'(Z')})\quoi,
$$
qui r\'esulte de l'hypoth\`ese que $f'$ est un hom\'eomorphisme
dans.

N.B. Dans ce dernier raisonnement, supposant d\'ej\`a que~$f$ est
quasi-compact, on n'a pas utilis\'e que~$g$ est quasi-compact, mais
seulement que
$g$
est fid\`element plat. Donc c'est sous cette hypoth\`ese qu'on
peut descendre la propri\'et\'e \og hom\'eomorphisme dans\fg, ou
\og hom\'eomorphisme sur\fg, ou encore gr\^ace au raisonnement
pr\'ec\'edent, la propri\'et\'e \og $f'$ est quasi-compact et
fait de~$f'(X')$ un espace topologique quotient de~$X'$\fg.

Nous dirons qu'un morphisme~$f\colon X\to Y$ de pr\'esch\'emas est
\emph{universellement ouvert} (\resp \emph{universellement fermé},
\resp \emph{universellement bicontinu}, etc.)
\index{universellement ouvert (\resp universellement ferm\'e, \resp
universellement bicontinu, etc.)|hyperpage}si pour tout changement de base~$Y'\to Y$, $f'\colon X'\to Y'$ est un
morphisme ouvert (\resp ferm\'e, \resp un
hom\'eomorphisme
sur l'espace image). On tire alors de~\ref{VIII.4.6}:

\begin{corollaire}
\label{VIII.4.7}
Sous
les conditions de~\ref{VIII.4.6}, pour que~$f$ soit
universellement ouvert, (\resp universellement ferm\'e, \resp un
hom\'eomorphisme dans universel, \resp un hom\'eomorphisme
universel), il faut et il suffit que $f'$ le soit.
\end{corollaire}

\begin{corollaire}
\label{VIII.4.8}
Sous les conditions de~\ref{VIII.4.6}, pour que~$f$ soit
s\'epar\'e (\resp propre) il faut et il suffit que $f'$ le soit.
\end{corollaire}

Dire que~$f$ est s\'epar\'e signifie que le morphisme
diagonal~$X\to X\times_{Y}X$ est ferm\'e ou aussi universellement
ferm\'e, et la premi\`ere assertion
de
\ref{VIII.4.8} r\'esulte donc de~\ref{VIII.4.7}. Dire que~$f$ est
propre signifie que $f$ satisfait les conditions a) $f$ est de type
fini b) $f$ est s\'epar\'e c) $f$ est universellement
ferm\'e. La condition a) se descend bien en vertu de~\ref{VIII.3.3},
b) aussi d'apr\`es ce qu'on vient de voir, enfin c) \'egalement
par~\ref{VIII.4.7}.

\begin{remarques}
\label{VIII.4.9}
Rappelons que lorsque $g\colon Y'\to Y$ est un morphisme plat de type
fini, avec~$Y$ localement noeth\'erien, alors $g$ est un morphisme
ouvert (VI~\ref{IV.6.6}) ce qui est un r\'esultat plus
pr\'ecis que~\ref{VIII.4.3}. On notera cependant que si~$f$ est un
morphisme fid\`element plat et quasi-compact de pr\'esch\'emas
noeth\'eriens, alors $f$ n'est pas en g\'en\'eral un morphisme
ouvert. Soit par exemple $Y$ un sch\'ema irr\'eductible dont le
point g\'en\'erique~$y$ n'est pas ouvert (par exemple une courbe
alg\'ebrique), et prenons pour~$Y'$ le sch\'ema somme~$Y\amalg
\Spec(\kres(y))$, alors l'image par le morphisme structural~$Y'\to Y$ de
la partie ouverte~$\Spec(\kres(y))$ n'est pas une partie ouverte
de~$Y$. Le lecteur remarquera \'egalement que divers \'enonc\'es
du pr\'esent expos\'e deviennent faux si on y abandonne
l'hypoth\`ese que le morphisme fid\`element plat envisag\'e est
aussi quasi-compact, le cas type mettant les \'enonc\'es en
d\'efaut \'etant celui o\`u on prend pour~$Y'$ le sch\'ema
somme des spectres des anneaux locaux des points de~$Y$. Par exemple,
prenant encore pour~$Y$ une courbe alg\'ebrique irr\'eductible, et
pour~$Z$ la partie de~$Y$ r\'eduite au point g\'en\'erique, son
image inverse dans~$Y'$ est ouverte, sans que~$Z$ soit ouverte.
\end{remarques}

\subsection{}
\label{VIII.4.10}
Divers
\'enonc\'es du pr\'esent expos\'e restent valables en y
rempla\c cant l'hypoth\`ese que~$Y'$ soit plat sur~$Y$ par la
suivante: il existe un Module de type fini~$F$ sur~$Y'$, de
support~$Y'$, plat par rapport \`a~$Y$; l'hypoth\`ese de
fid\`ele platitude sera alors remplac\'ee par la
pr\'ec\'edente, plus l'hypoth\`ese que $Y'\to Y$ est
surjectif. Ceci s'applique aux deux premi\`eres assertions
dans~\ref{VIII.1.10}, \`a~\ref{VIII.3.3}, \ref{VIII.3.5},
\ref{VIII.4.1} et par suite \`a tous les r\'esultats du
pr\'esent num\'ero.

\section{Descente de morphismes de pr\'esch\'emas}
\label{VIII.5}

\begin{proposition}
\label{VIII.5.1}
Soit $g\colon S'\to S$ un morphisme de pr\'esch\'emas.
\begin{enumerate}
\item[a)] Supposons que~$g$ soit surjectif, et que l'homomorphisme
$$
g^*\colon \mathcal{O}_{S}\to g_*(\mathcal{O}_{S'})
$$
soit injectif, alors~$g$ est un \'epimorphisme dans la cat\'egorie
des pr\'esch\'emas, et m\^eme dans la cat\'egorie des espaces
annel\'es.
\item[b)] Supposons que~$g$ soit surjectif et fasse de~$S$ un espace
topologique quotient de~$S'$. Soit~$S''=S'\times_{S}S'$ et
soit~$h\colon S''\to S$ le morphisme structural, consid\'erons le
diagramme d'homomorphismes canonique:
$$
\xymatrix@C=.5cm{{\mathcal{O}_{S}}\ar[r]&{g_*(\mathcal{O}_{S'})}\ar@<2pt>[r]\ar@<-2pt>[r]&{h_*(\mathcal{O}_{S''})}}\quoi,
$$
et supposons ce diagramme \emph{exact}. Alors $g$ est un
\'epimorphisme effectif
\index{epimorphisme effectif@\'epimorphisme effectif|hyperpage}dans la cat\'egorie des pr\'esch\'emas (et aussi dans la
cat\'egorie des espaces annel\'es), \ie le diagramme
$$
\xymatrix@C=.5cm{{S}&{S'}\ar[l]&{S''}\ar@<2pt>[l]\ar@<-2pt>[l]}
$$
est exact.
\end{enumerate}
\end{proposition}

\subsubsection*{D\'emonstration} a) Il faut montrer qu'un morphisme d'espaces
annel\'es~$f\colon S\to Z$ est connu quand on
conna\^it~$fg$. Or comme~$g$ est surjectif, on conna\^it
l'application ensembliste~$f_{0}$ sous-jacente \`a~$f$, reste \`a
d\'eterminer l'homomorphisme de faisceaux
d'anneaux~$\mathcal{O}_{Z}\to \mathcal{O}_{S}$, ou ce qui revient au
m\^eme l'homomorphisme
de faisceaux d'anneaux
$$
u\colon f_{0}^{-1}(\mathcal{O}_{Z})\to \mathcal{O}_{S}
$$
d\'efini par~$f$. On conna\^it d\'ej\`a l'homomorphisme
$$
(fg)_{0}^{-1}(\mathcal{O}_{Z})=
g_{0}^{-1}(f_{0}^{-1}(\mathcal{O}_{Z}))\to \mathcal{O}_{S'}
$$
d\'efini par~$fg$, ou ce qui revient au m\^eme, on dispose d'un
homomorphisme
$$
f_{0}^{-1}(\mathcal{O}_{Z})\to
g_{0*}(\mathcal{O}_{S'})=g_*(\mathcal{O}_{S'})\quoi.
$$
On constate aussit\^ot que ce dernier n'est autre que le compos\'e
de~$g^*\colon \mathcal{O}_{S}\to g_*(\mathcal{O}_{S'})$ et
de~$u$, et comme $g^*$ est injectif, $u$ est connu quand on
conna\^it~$g^*u$. [N.B. on n'a pas utilis\'e \'evidemment
que~$g\colon S'\to S$ est un morphisme de pr\'esch\'emas,
l'\'enonc\'e vaudrait pour un morphisme quelconque d'espaces
annel\'es; la m\^eme remarque vaut pour b), tant dans la
cat\'egorie des espaces annel\'es, que dans la cat\'egorie des
espaces annel\'es en anneaux locaux. Noter aussi que si~$g$ est un
morphisme de pr\'esch\'emas pas n\'ecessairement surjectif, mais
tel que~$g^*\colon \mathcal{O}_{S}\to g_*(\mathcal{O}_{S'})$
soit injectif, alors pour deux morphismes~$f_{1},f_{2}$ de~$S$ dans un
\emph{schéma}~$Z$ tel que~$f_{1}g=f_{2}g$, on a $f_{1}=f_{2}$; en
effet, si~$I$ est l'Id\'eal sur~$S$ qui d\'efinit le
sous-pr\'esch\'ema de~$S$ des co\"incidences de~$f_{1},f_{2}$
(image inverse du sous-pr\'esch\'ema diagonal de~$Z\times Z$
par~$(f_{1},f_{2})$), on voit que~$I$ est contenu
dans~$\Ker(g^*)$].

b) On doit montrer que pour tout espace annel\'e~$Z$, le diagramme
suivant d'applications
$$
\xymatrix@C=.5cm{{\mathrm{Hom}(S,Z)}\ar[r]&{\mathrm{Hom}(S',Z)}\ar@<2pt>[r]\ar@<-2pt>[r]&{\mathrm{Hom}(S'',Z)}}
$$
est exact, et qu'il en est de m\^eme lorsque~$Z$ est un espace
annel\'e en anneaux locaux et qu'on se borne aux homomorphismes
d'espaces annel\'es en anneaux locaux. Comme on sait d\'ej\`a par a)
que la premi\`ere application est injective, il reste \`a voir que
si~$f'\colon S'\to Z$ est un homomorphisme d'espaces annel\'es tel
que~$f'p_{1}=f'p_{2}$,
alors~$f'$ est de la forme~$fg$, o\`u~$f\colon S\to Z$ est un
homomorphisme d'espaces annel\'es. Comme~$g$ est surjectif, il est
alors \'evident que si~$f'$ est un morphisme d'espaces annel\'es
en anneaux locaux, il en sera de m\^eme pour~$f$.

De l'hypoth\`ese sur~$f'$ r\'esulte que l'application ensembliste
sous-jacente~$f_{0}'$ est constante sur les fibres de
l'application~$g_{0}$, donc comme cette derni\`ere est surjective,
$f_{0}'$ se factorise de fa\c con unique en~$f_{0}'=f_{0}g_{0}$,
o\`u~$f_{0}\colon S\to Z$ est une application, n\'ecessairement
continue puisque~$g_{0}$ identifie~$S$ \`a un espace topologique
quotient de~$S'$. Consid\'erons maintenant l'homomorphisme
$$
f_{0}^{-1}(\mathcal{O}_{Z})\to g_*(\mathcal{O}_{S'})
$$
d\'eduit de l'homomorphisme~$(f_{0}g_{0})^{-1}(\mathcal{O}_{Z})\to
\mathcal{O}_{S'}$ correspondant
\`a~$f'$. L'hypoth\`ese~$f'p_{1}=f'p_{2}$ s'interpr\`ete alors
en disant que les compos\'es de l'homomorphisme pr\'ec\'edent
avec les deux homomorphismes
$\xymatrix@C=.5cm{{g_*(\mathcal{O}_{S'})}\ar@<2pt>[r]\ar@<-2pt>[r]&{h_*(\mathcal{O}_{S''})}}$ sont les m\^emes donc d'apr\`es
l'hypoth\`ese b) il se factorise par un morphisme
$$
f_{0}^{-1}(\mathcal{O}_{Z})\to \mathcal{O}_{S}\quoi.
$$
Ce dernier d\'efinit un morphisme d'espaces annel\'es~$f\colon
S\to Z$, qui est le morphisme cherch\'e.

\begin{theoreme}
\label{VIII.5.2}
Soit $\cal{F}$ la cat\'egorie fibr\'ee des fl\`eches dans la
cat\'egorie~$\Sch$ des pr\'esch\'emas
\textup{(VI~\ref{VI.11}.a)}. Alors tout morphisme~$g\colon S'\to S$
fid\`element plat et quasi-compact est un morphisme
de~$\cal{F}$-descente (ou encore, comme on dit, un morphisme de
descente dans~$\Sch$).
\index{descente (morphisme de)|hyperpage}\end{theoreme}

Cela signifie donc ceci: soit~$S''=S'\times_{S}S'$, et pour deux
pr\'esch\'emas~$X,Y$ sur~$S$, consid\'erons leurs images
inverses~$X',Y'$ sur~$S'$ et leurs images inverses~$X'',Y''$
sur~$S''$, d'o\`u un diagramme d'applications
$$
\xymatrix@C=.5cm{{\mathrm{Hom}_{S}(X,Y)}\ar[r]&{\mathrm{Hom}_{S'}(X',Y')}\ar@<2pt>[r]\ar@<-2pt>[r]&{\mathrm{Hom}_{S''}(X'',Y'')}}\quoi;
$$
ces notations pos\'ees, \ref{VIII.5.2} signifie que ce diagramme est
exact. On notera qu'il
n'est pas vrai en g\'en\'eral que~$g$ soit un morphisme de
descente effective, \ie que pour tout pr\'esch\'ema~$X'$
sur~$S'$, toute donn\'ee de descente sur~$X'$ relativement
\`a~$g\colon S'\to S$ soit effective. La question de
l'effectivit\'e, souvent d\'elicate, sera examin\'ee au
\No \ref{VIII.7}.

On a vu \cite{VIII.D}, (compte tenu que dans~$\Sch$ les produits fibr\'es
existent) que l'\'enonc\'e~\ref{VIII.5.2} \'equivaut au suivant:

\begin{corollaire}
\label{VIII.5.3}
Un morphisme fid\`element plat et quasi-compact de
pr\'esch\'emas est un \'epi\-mor\-phisme effectif universel.
\end{corollaire}

Comme un morphisme fid\`element plat et quasi-compact reste tel par
toute extension de la base, on est ramen\'e \`a prouver que c'est
un \'epimorphisme effectif. On applique alors le
crit\`ere~\ref{VIII.5.1} b), qui donne le r\'esultat voulu, compte
tenu de~\ref{VIII.4.3} et~\ref{VIII.1.7}.

\begin{corollaire}
\label{VIII.5.4}
Soit $g\colon S'\to S$ un morphisme fid\`element plat et
quasi-compact, $f\colon X\to Y$ un~$S$-morphisme, $f'\colon X'\to Y'$
le~$S'$-morphisme qui s'en d\'eduit par le changement de base~$S'\to
S$. Pour que~$f$ soit un isomorphisme, il faut et il suffit que~$f'$
le soit.
\end{corollaire}

En effet, si $f'$ est un isomorphisme, c'est aussi un isomorphisme
pour les structures de descente naturelles sur~$X',Y'$, et comme le
foncteur~$X\mto X'$ de~$\Sch_{/S}$ dans la cat\'egorie des
pr\'esch\'emas sur~$S'$ munis d'une donn\'ee de descente
relativement \`a~$g$ est pleinement fid\`ele par~\ref{VIII.5.2},
la conclusion voulue appara\^it.

\begin{corollaire}
\label{VIII.5.5}
Sous les conditions de~\ref{VIII.5.4}, pour que~$f$ soit une immersion
ferm\'ee (\resp une immersion ouverte, \resp une immersion
quasi-compacte) il faut et suffit que~$f'$ le soit.
\end{corollaire}

On peut supposer comme d'habitude~$Y=S$, et il y a \`a prouver
seulement le \og il suffit\fg. Notons que le fait que~$X'/Y'$ soit muni
d'une donn\'ee de descente relativement \`a~$g\colon Y'\to Y$, et
que le morphisme structural~$f'\colon X'\to Y'$ soit une immersion
donc un monomorphisme, implique que les deux sous-objets de~$Y''$
images
inverses de~$X'/Y'$ par l'une et l'autre projection de~$S''$ dans~$S'$
sont les m\^emes. Si~$f'$ est une immersion ferm\'ee, il en
r\'esulte en vertu de~\ref{VIII.1.9} qu'il existe un
sous-pr\'esch\'ema ferm\'e $X_{1}$ de~$Y$ dont l'image inverse
par~$g\colon Y'\to Y$ est~$X'$. Donc par l'unicit\'e de la solution
d'un probl\`eme de descente relativement \`a un morphisme de
$\cal{F}$-descente,
il
r\'esulte que~$X_{1}$ est $Y$-isomorphe \`a~$X$, donc $f\colon
X\to Y$ est une immersion ferm\'ee. On proc\`ede de m\^eme pour
une immersion ouverte, en utilisant~\ref{VIII.4.4}. Si enfin $f'$ est
une immersion quasi-compacte, $f$ est quasi-compact en vertu
de~\ref{VIII.3.3}, donc on peut appliquer \`a la partie~$f(X)$
de~$Y$ le crit\`ere~\ref{VIII.4.5}, qui prouve que~$f(X)$ est
localement ferm\'e puisque son image inverse~$f'(X')$ dans~$Y'$
l'est. Rempla\c cant alors~$Y$ par une partie ouverte dans
laquelle~$f(X)$ soit ferm\'ee, on est ramen\'e au cas o\`u~$f'$
est une immersion ferm\'ee, donc $f$ l'est en vertu de ce qui
pr\'ec\`ede.

\begin{corollaire}
\label{VIII.5.6}
Sous les conditions de~\ref{VIII.5.4}, pour que~$f$ soit affine, il
faut et il suffit que~$f'$ le soit.
\end{corollaire}

On proc\`ede comme dans~\ref{VIII.5.5}, en utilisant~\ref{VIII.2.1}
(On peut aussi utiliser le crit\`ere cohomologique de Serre [EGA II
5.2], qui d\'emontre~\ref{VIII.5.6} sans utiliser de technique de
descente).

\begin{corollaire}
\label{VIII.5.7}
Sous les conditions de~\ref{VIII.5.4}, pour que~$f$ soit entier
(\resp fini, \resp fini et localement libre) il faut et il suffit
que~$f'$ le soit.
\end{corollaire}

Il y a \`a prouver seulement le \og il suffit\fg, et comme d'habitude
on peut supposer $Y=S$, $Y$ affine, et $Y'$ affine. Comme
l'hypoth\`ese implique que~$f'$ est affine, il en est de m\^eme
de~$f$ d'apr\`es~\ref{VIII.5.6}, donc~$X$ et par suite~$X'$ sont
affines. Soient~$A,A',B,B'=B\otimes_{A}A'$ les anneaux
de~$Y,Y',X,X'$. On a $B=\varinjlim_i B_{i}$, o\`u~$B_{i}$
parcourt les sous-$A$-alg\`ebres de~$B$ qui sont de type fini
sur~$A$, d'o\`u $B'=\varinjlim_i B_{i}'$, o\`u
les~$B_{i}'$ sont des sous-alg\`ebres de type fini de la
$A'$-alg\`ebre~$B'$. Si $B'$ est entier
sur~$A'$,
les~$B_{i}'$ sont des modules de type fini sur~$A'$, donc $A'$
\'etant fid\`element plat sur~$A$, les~$B_{i}$ sont des modules de
type fini sur~$A$, \ie $B$ est entier sur~$A$. On voit de m\^eme
que si~$B'$ est fini sur~$A'$, $B$ l'est sur~$A$. M\^eme conclusion
pour \og localement libre de type fini\fg, \cf \ref{VIII.1.11}.

\begin{corollaire}
\label{VIII.5.8}
Sous
les conditions de~\ref{VIII.5.4}, supposons~$f$ quasi-compact et
soient~$\mathcal{L}$ un Module inversible sur~$X$, et~$\mathcal{L}'$
son image inverse sur~$X'$. Pour que~$\mathcal{L}$ soit ample
(\resp tr\`es ample) relativement \`a~$f$, il faut et il suffit
que~$\cal{L}'$
soit ample (\resp tr\`es ample) relativement \`a~$f'$.
\end{corollaire}

Il y a \`a prouver seulement le \og il suffit\fg. L'hypoth\`ese
sur~$\mathcal{L}$ implique en tous cas que~$f'$ est s\'epar\'e,
donc $f$ est s\'epar\'e par~\ref{VIII.4.8}, et comme~$f$ est
quasi-compact, et~$g\colon Y'\to Y$ est plat, le calcul des images
directes par recouvrements affines montre qu'on a des isomorphismes
$$
g^*(f_*(\mathcal{L}^{\otimes n}))\isomto
f_*'(\mathcal{L}'^{\otimes n})
$$
pour tout entier~$n$, donc on a un isomorphisme
$$
g^*(\mathcal{S})\isomto \mathcal{S}'\quoi,
$$
o\`u~$\mathcal{S}$ (\resp $\mathcal{S}'$) d\'esigne l'Alg\`ebre
gradu\'ee quasi-coh\'erente sur~$Y$ (\resp sur~$Y'$) somme directe
des~$f_*(\mathcal{L}^{\otimes n})$
(\resp des~$f_*'(\mathcal{L}'^{\otimes n})$) pour~$n\geq
0$. Notons que pour tout~$n\geq 0$, le conoyau de l'homomorphisme
canonique~$f'^*(\mathcal{S}_{n}')\to \mathcal{L}'^{\otimes n}$
est l'image inverse par~$X'\to X$ du conoyau
de~$f^*(\mathcal{S}_{n})\to \mathcal{L}^{\otimes n}$, donc son
support~$Z_{n}'$ est l'image inverse du support~$Z_{n}$. Si
$\cal{L}'$
est ample l'intersection des~$Z_{n}'$ est vide, donc comme~$X'\to X$
est surjectif, l'intersection des~$Z_{n}$ est vide, \ie on a un
morphisme canonique
$$
j\colon X\to \Proj(\mathcal{S})
$$
(EGA~II 3). D'ailleurs, le morphisme analogue
$$
j'\colon X'\to \Proj(\mathcal{S}')
$$
n'est autre que celui qui est d\'eduit du pr\'ec\'edent par le
changement de base~$Y'\to Y$ (\loccit). Ceci pos\'e, dire que
$\mathcal{L}'$ est ample relativement \`a~$f'$ signifie que~$j'$ est
une immersion, d'ailleurs n\'ecessairement quasi-compacte
puisque~$f'$ est quasi-compact. Donc en vertu de~\ref{VIII.5.5} $j$
est une immersion, \ie $\mathcal{L}$ est ample relativement
\`a~$f$. --- On proc\`ede de fa\c con toute analogue dans le cas
de \og tr\`es
ample\fg, en se bornant ci-dessus \`a~$n=1$, et en rempla\c cant la
consid\'eration de~$\Proj(\mathcal{S})$ par celle du fibr\'e
projectif~$\mathcal{P}(\mathcal{S}_{1})$ associ\'e
\`a~$\mathcal{S}_{1}$.

Rappelons (EGA~II 5.1.1) qu'un morphisme quasi-compact~$f$ est dit
\emph{quasi-affine} si pour tout ouvert affine~$U$ dans~$Y$,
$f^{-1}(U)$ est un pr\'esch\'ema isomorphe \`a un
sous-sch\'ema ouvert d'un sch\'ema affine. On montre (\loccit)
qu'il revient au m\^eme de dire que~$\mathcal{O}_{X}$ est ample (ou
aussi: tr\`es ample) relativement \`a~$f$. Donc~\ref{VIII.5.8}
implique:

\begin{corollaire}
\label{VIII.5.9}
Sous les conditions de~\ref{VIII.5.4}, et supposant~$f$ quasi-compact,
pour que~$f$ soit quasi-affine, il faut et il suffit que~$f'$ le soit.
\end{corollaire}

\begin{remarques}
\label{VIII.5.10}
L'exemple de vari\'et\'e non projective de Hironaka montre qu'on
peut avoir un morphisme propre~$f\colon X\to Y$ de vari\'et\'es
alg\'ebriques non singuli\`eres (avec~$Y$ projective), tel que~$Y$
soit r\'eunion de deux ouverts~$Y_{i}$ tels
que~$X_{i}=X\times_{Y}Y_{i}$ soit projectif sur~$Y_{i}$, mais~$f$
n'\'etant pas projectif. Donc posant~$Y'=Y_{1}\amalg Y_{2}$, $Y'$
est fid\`element plat et quasi-compact (et m\^eme quasi-fini)
sur~$Y$, $f'\colon X'\to Y'$ est projectif, mais~$f$ n'est pas
projectif. Il faut donc faire attention que pour
appliquer~\ref{VIII.5.8}, et d\'eduire du fait que~$f'$ est
projectif la m\^eme conclusion sur~$f$, il faut disposer d\'ej\`a
sur~$X'$ d'un Module inversible~$\mathcal{L}'$ ample pour~$f'$,
\emph{muni d'une donnée de descente relativement à}~$X'\to X$,
(ce qui permet de consid\'erer~$\mathcal{L}'$ comme l'image inverse
d'un Module inversible~$\mathcal{L}$ sur~$X$, qui sera alors ample
pour~$f$ gr\^ace \`a~\ref{VIII.5.8}). Lorsque~$g\colon S'\to S$
est fini et localement libre, voir cependant~\ref{VIII.7.7}.
\end{remarques}

\section[Application aux morphismes finis et quasi-finis]{Application aux morphismes finis et quasi-finis\footnotemark}
\label{VIII.6}
\footnotetext{\Cf EGA IV 18.12 pour des g\'en\'eralisations \`a des pr\'esch\'emas non n\'ecessairement localement noeth\'eriens}

Nous allons d\'emontrer les deux th\'eor\`emes suivants:

\begin{theoreme}
\label{VIII.6.1}
Soit~$f\colon X\to Y$ un morphisme \emph{propre à fibres finies},
avec~$Y$ localement noeth\'erien. Alors~$f$ est fini.
\end{theoreme}

\begin{theoreme}
\label{VIII.6.2}
Soit~$f\colon X\to Y$ un morphisme \emph{quasi-fini} et
\emph{séparé}, avec~$Y$ localement noeth\'erien. Alors~$f$
est quasi-affine, et a fortiori quasi-projectif.
\end{theoreme}

\begin{remarques}
\label{VIII.6.3}
Le
th\'eor\`eme~\ref{VIII.6.1} est bien connu, et d\^u \`a
Chevalley dans le cas de vari\'et\'es alg\'ebriques; on en
trouvera aussi une d\'emonstration simple dans [EGA III 4],
utilisant le \og th\'eor\`eme des fonctions holomorphes\fg. La
d\'emonstration donn\'ee ici n'utilise pas ce dernier
th\'eor\`eme, mais par contre la th\'eorie de descente; nous
la donnons comme prime au lecteur, car on l'a \og pour rien\fg en
m\^eme temps que celle de~\ref{VIII.6.2}. Rappelons aussi ([EGA III
4] ou \cite{VIII.1}) que la forme globale du \og Main Theorem\fg de Zariski,
d\'eduit du \og th\'eor\`eme des fonctions holomorphes\fg, affirme
que si~$f\colon X\to Y$ est quasi-fini et \emph{quasi-projectif}, $Y$
\'etant noeth\'erien, alors~$X$ est $Y$-isomorphe \`a un
sous-pr\'esch\'ema ouvert d'un~$Y$-pr\'esch\'ema
\emph{fini}~$Z$. La conjonction du \og Main Theorem\fg et
de~\ref{VIII.6.2} s'\'enonce donc ainsi:
\end{remarques}

\begin{corollaire}
\label{VIII.6.4}
Soit~$f\colon X\to Y$ un morphisme quasi-fini et s\'epar\'e,
avec~$Y$ noeth\'erien. Alors~$X$ est $Y$-isomorphe \`a un
sous-pr\'esch\'ema ouvert d'un~$Y$-pr\'esch\'ema fini~$Z$.
\end{corollaire}

Une autre cons\'equence int\'eressante de~\ref{VIII.6.2} pour la
th\'eorie de la descente sera donn\'ee avec~\ref{VIII.7.9}.

\subsubsection*{D\'emonstration de~\ref{VIII.6.1} et~\ref{VIII.6.2}}
Nous admettrons le fait suivant, dont la d\'emonstration est
facile\footnote{\Cf EGA IV 8}:

\begin{lemme}
\label{VIII.6.5}
Soit~$X$ un pr\'esch\'ema de type fini sur~$Y$ localement
noeth\'erien, et soit~$y\in Y$. Pour qu'il existe un voisinage
ouvert~$U$ de~$y$ tel que~$X|U$ soit fini (\resp quasi-affine,
\resp\ ...) sur~$U$, il faut et il suffit
que~$X\times_{Y}\Spec(\mathcal{O}_{y})$ soit fini (\resp quasi-affine,
\resp\ ...) sur~$\Spec(\mathcal{O}_{y})$.
\end{lemme}

Comme d'autre part la propri\'et\'e pour~$f\colon X\to Y$
d'\^etre fini, \resp quasi-affine, est locale sur~$Y$, on est
ramen\'e pour prouver~\ref{VIII.6.1} et~\ref{VIII.6.2} au cas
o\`u~$Y$ est le spectre d'un anneau local, et est a fortiori de
dimension finie. (N.B. on appelle \emph{dimension d'un
préschéma}~$Y$ le sup des dimensions de Krull de ses anneaux
locaux). Nous proc\'edons par r\'ecurrence sur~$$
n=\dim(Y)\quoi,
$$
l'assertion
\'etant triviale pour~$n<0$. Nous pouvons donc supposer~$n\geq 0$,
et l'assertion d\'emontr\'ee pour les dimensions~$n'<n$. On peut
\`a nouveau supposer que~$Y$ est le spectre d'un anneau local
noeth\'erien~$A$, de dimension~$n$. Notons que les hypoth\`eses
faites dans~\ref{VIII.6.1} et~\ref{VIII.6.2} sont stables par
changement de base (on s'en est d\'ej\`a servi dans la r\'eduction
du d\'ebut), elles resteront vraies apr\`es le changement de
base~$\Spec(\widehat{A})\to \Spec(A)$. Comme ce dernier est
fid\`element plat et quasi-compact, les
\'enonc\'es~\ref{VIII.5.7} et~\ref{VIII.5.9} nous ram\`enent au
cas o\`u~$A$ est de plus complet. Utilisant alors le fait que tout
anneau local noeth\'erien~$B$ sur~$A$ quasi-fini sur~$A$ est fini
sur~$A$, et le fait que~$X$ est s\'epar\'e sur~$Y$ et la fibre
de~$y$ est form\'ee de points isol\'es, on trouve une
d\'ecomposition
$$
X=X'\amalg X''
$$
o\`u~$X'$ est \emph{fini} sur~$Y$, et o\`u la fibre de~$X''$
en~$y$ est vide. Si~$X$ est propre sur~$Y$, il en est de m\^eme
de~$X''$, donc son image dans~$Y$ est ferm\'ee, et comme elle ne
contient pas~$y$, elle est vide, donc~$X''=\emptyset$ donc~$X=X'$,
ce qui montre que~$X$ est fini sur~$Y$ et d\'emontre~\ref{VIII.6.1}
(N.B. l'hypoth\`ese de r\'ecurrence est ici inutile). Si~$X$ est
quasi-fini sur~$Y$, $X''$ l'est aussi, or~$X''$ se trouve en fait sur
l'ouvert~$Y-\{ y\}$
de~$Y$, \emph{qui est de dimension} $<n$. En vertu de l'hypoth\`ese
de r\'ecurrence, $X''$ est quasi-affine
sur~$Y-\{ y\}$,
donc aussi sur~$Y$, il en est \'evidemment de m\^eme de~$X'$, donc
aussi de leur somme~$X$, ce qui prouve~\ref{VIII.6.2}.

\begin{remarque}
\label{VIII.6.6}
Les th\'eor\`emes~\ref{VIII.6.1} et~\ref{VIII.6.2} restent
valables si on ne suppose plus~$Y$ localement noeth\'erien, \`a
condition de sp\'ecifier que l'on suppose~$f$ de pr\'esentation
finie
(\cf \ref{VIII.3.6}).
En effet, on peut encore supposer~$Y$ affine, et alors on v\'erifie
sans difficult\'e que la situation~$f\colon X\to Y$ est
d\'eduite,
par un changement de base~$Y\to Y_{0}$, d'une situation~$f_{0}\colon
X_{0}\to Y_{0}$ satisfaisant les m\^emes hypoth\`eses que~$f$,
avec~$Y_{0}$ \emph{noethérien}. Donc d'apr\`es le
r\'esultat~\ref{VIII.6.1} \resp \ref{VIII.6.2}, $f_{0}$ est fini
\resp quasi-affine, donc il en est de m\^eme de~$f$. Ce genre de
raisonnement est souvent utile pour se d\'ebarrasser
d'hypoth\`eses noeth\'eriennes,
(qui finissent toujours par \^etre g\^enantes dans les
applications).
\end{remarque}

\section{Crit\`eres d'effectivit\'e pour une donn\'ee de descente}
\label{VIII.7}

Consid\'erons comme d'habitude un morphisme de pr\'esch\'emas
$$
g\colon S'\to S
$$
et un $S'$-pr\'esch\'ema~$X'$. Conform\'ement aux faits
g\'en\'eraux~VII,~9\footnote{n'existe pas; voir expos\'e 212 du
S\'eminaire Bourbaki (1962)},
la donn\'ee d'une donn\'ee de descente sur~$X'$, relativement
\`a~$g$, est \'equivalente \`a la donn\'ee d'un couple
d'\'equivalence \cite{VIII.3}:
$$
\xymatrix@C=.5cm{{q_{1},q_{2}\colon X''}\ar@<2pt>[r]\ar@<-2pt>[r]&{X'}}
$$
tel que le morphisme structural~$X'\to S'$ soit compatible avec ce
couple et le couple d'\'equivalence
$$
\xymatrix@C=.5cm{{p_{1},p_{2}\colon
S''=S'\times_{S}S'}\ar@<2pt>[r]\ar@<-2pt>[r]&{S'}}
$$
d\'efini par~$g$, et tel que les deux carr\'es (ou l'un des deux,
cela revient au m\^eme par raison de sym\'etrie) extraits du
diagramme correspondant
$$
\xymatrix{{X'}\ar[d]&{X''}\ar@<2pt>[l]\ar@<-2pt>[l]\ar[d]\\{S'}&{S''}\ar@<2pt>[l]\ar@<-2pt>[l]}
$$
(en utilisant soit~$p_{1},q_{1}$, soit~$p_{2},q_{2}$), soit
\emph{cartésien}. Une solution du probl\`eme de descente
pos\'e par cette donn\'ee de descente, \ie un objet~$X$ sur~$S$
muni d'un isomorphisme~$X\times_{S}S'\isomfrom X'$ compatible avec les
donn\'ees de descente, \'equivaut \`a la donn\'ee d'un
carr\'e \emph{cartésien}
$$
\xymatrix{{X}\ar[d]&{X'}\ar[l]_{h}\ar[d]\\{S}&{S'}\ar[l]_{g}}
$$
satisfaisant~$hq_{1}=hq_{2}$.

Comme
l'ensemble des morphismes fid\`element plats et quasi-compacts est
stable par changement de base, et qu'un morphisme fid\`element plat
et quasi-compact est un \'epimorphisme effectif en vertu
de~\ref{VIII.5.3}, il r\'esulte de la th\'eorie g\'en\'erale
\cite{VIII.D}:

\begin{proposition}
\label{VIII.7.1}
Supposons~$g\colon S'\to S$ fid\`element plat et quasi-compact. Pour
qu'une don\-n\'ee de descente sur~$X'$ relativement \`a~$g$ soit
effective, il faut et il suffit que la relation
d'\'equivalence~$R=(q_{1},q_{2})$ qu'elle d\'efinit soit effective
(\ie le quotient~$X'/R$ existe et $X''$ devient le carr\'e
fibr\'e de~$X'$ sur~$X'/R$), et que le morphisme canonique~$X'\to
X'/R$ soit fid\`element plat et quasi-compact.
\end{proposition}

Ainsi la question de l'effectivit\'e d'une donn\'ee de descente
est un cas particulier de la question d'effectivit\'e d'un graphe
d'\'equivalence, et divers crit\`eres d'effectivit\'e donn\'es
dans ce num\'ero peuvent s'obtenir de cette fa\c con. N\'eanmoins on dispose dans le contexte de la descente du
th\'eor\`eme~\ref{VIII.2.1}, qui implique que \emph{si $X'$ est
affine sur~$S'$, toute donnée de descente sur~$X'$ relativement
à~$g$ est effective}, \'enonc\'e qui n'a pas d'analogue pour
le passage au quotient par un graphe d'\'equivalence plat
g\'en\'eral. Tous les crit\`eres d'effectivit\'e que nous
donnons ici peuvent aussi \^etre consid\'er\'es comme
d\'eduits du pr\'ec\'edent.

Soit~$U'$ un sous-pr\'esch\'ema de~$X'$ (ou plus
g\'en\'eralement un sous-objet de~$X'$ dans la
cat\'egorie~$\Sch$); on dit que~$U'$ est \emph{stable par la
donnée de descente} sur~$X'$, si on peut trouver sur~$U'$ une
donn\'ee de descente relativement \`a~$g$, telle que
l'immersion~$U'\to X'$ soit compatible avec les donn\'ees de
descente. Cela signifie aussi que les images inverses de~$U'$
dans~$X''$ par~$q_{1}$ et~$q_{2}$ sont les m\^emes (ou aussi, comme
on dit, que~$U'$ est \emph{stable par la relation
d'équivalence}~$R$), et bien entendu la donn\'ee de descente en
question sur~$U'$ est alors unique, et dite \emph{donnée de
descente induite}
\index{donn\'ee de descente induite|hyperpage}par celle de~$X'$. Ceci pos\'e:

\begin{proposition}
\label{VIII.7.2}
Soit~$(X_{i}')$ un recouvrement de~$X'$ par des ouverts~$X_{i}'$
stables
par la donn\'ee de descente. Pour que la donn\'ee de descente
sur~$X'$ soit effective, il faut et il suffit qu'il en soit de
m\^eme des donn\'ees de descente induites dans les~$X_{i}'$.
\end{proposition}

C'est l\`a une cons\'equence facile de~\ref{VIII.7.1} par exemple,
et le d\'etail de la d\'emonstration est laiss\'e au lecteur.

\begin{corollaire}
\label{VIII.7.3}
Soit~$(S_{i})$ un recouvrement ouvert de~$S$, et pour tout~$i$
soient~$S_{i}'$ et~$X_{i}'$ d\'eduits de~$S'$
et
$X'$ par le changement de base~$S_{i}\to S$. Pour que la donn\'ee de
descente sur~$X'$ soit effective, il faut et il suffit que, pour
tout~$i$, la donn\'ee de descente sur~$X_{i}'$ relativement
\`a~$g_{i}\colon S_{i}'\to S_{i}$ soit effective.
\end{corollaire}

Ce crit\`ere nous ram\`ene toujours pratiquement au cas o\`u~$S$
est affine. Dans le cas o\`u~$S'$ est \'egalement affine, ce qui
est le cas le plus fr\'equent dans les applications, on~a:

\begin{corollaire}
\label{VIII.7.4}
Supposons~$S$ et~$S'$ affines. Pour que la donn\'ee de descente
sur~$X'$ soit effective, il faut et il suffit que~$X'$ soit
r\'eunion d'ouverts~$X_{i}'$ affines et stables par la donn\'ee de
descente.
\end{corollaire}

La suffisance provient de~\ref{VIII.7.2} et du fait que si~$X_{i}'$
est affine, il est affine sur~$S'$ et on peut
appliquer~\ref{VIII.2.1}. Pour la n\'ecessit\'e, on note que
si~$X'$ provient de~$X$, et si~$X$ est recouvert par des ouverts
affines~$X_{i}$, alors les~$X_{i}'=X_{i}\times_{S}S'$ sont des ouverts
affines stables par les donn\'ees de descente et recouvrant~$X'$.

\begin{corollaire}
\label{VIII.7.5}
Soit~$g\colon S'\to S$ un morphisme fid\`element plat, quasi-compact
et \emph{radiciel}. Alors~$g$ est un morphisme de descente
\emph{effective}, \ie pour tout~$X'$ sur~$S'$, toute donn\'ee de
descente sur~$X'$ relativement \`a~$g\colon S'\to S$ est effective.
\end{corollaire}

En effet, en vertu de~\ref{VIII.7.3} on peut supposer~$S$ affine, donc
comme~$S'$ est radiciel sur~$S$ donc s\'epar\'e, $S'$ est
s\'epar\'e. D'ailleurs pour tout~$x'\in X'$, la
fibre~$R(x')=q_{2}(q_{1}^{-1}(x'))$ de la relation d'\'equivalence
ensembliste d\'efinie
par la relation d'\'equivalence~$R$ est r\'eduite \`a un point,
car~$g$ \'etant radiciel, il en est de m\^eme de~$p_{1},p_{2}$ qui
s'en d\'eduisent par changement de base~$S'\to S$, donc aussi
de~$q_{1},q_{2}$ qui se d\'eduisent des pr\'ec\'edents par
changement de base~$X'\to S''$. Donc \emph{tout ouvert de~$X'$ est
stable} par la donn\'ee de descente. Recouvrons alors~$X'$ par des
ouverts affines~$X_{i}'$, ils sont affines sur~$S$ puisque~$S'$ est
s\'epar\'e, donc la donn\'ee de descente induite est effective
par~\ref{VIII.2.1}. On conclut alors par~\ref{VIII.7.2}.

On notera que~\ref{VIII.7.5} donne le seul cas connu d'un morphisme de
descente effective dans la cat\'egorie des pr\'esch\'emas, et
c'est probablement le seul cas en effet, m\^eme en se limitant aux
sch\'emas noeth\'eriens, ou aux sch\'emas de type fini sur un
corps.

Lorsqu'on suppose~$S$ localement noeth\'erien et~$S'$ de type fini
sur~$S$, l'\'enonc\'e~\ref{VIII.7.5} est aussi un cas particulier
du suivant (qui g\'en\'eralise la descente galoisienne de Weil et
la descente ins\'eparable de Cartier):

\begin{corollaire}
\label{VIII.7.6}
Soit~$g\colon S'\to S$ un morphisme fini localement libre
(\ie d\'efini par une Alg\`ebre sur~$S$ qui est un module
localement libre de type fini) et surjectif (donc~$g$ est
fid\`element plat et quasi-compact, donc un morphisme de
descente). Soit~$X'$ un~$S'$-pr\'esch\'ema muni d'une donn\'ee
de descente. Pour que cette donn\'ee soit effective, il faut et il
suffit que pour tout~$x'\in X'$, la
fibre~$R(x')=q_{2}(q_{1}^{-1}(x'))$ soit contenue dans un ouvert
affine (condition automatiquement v\'erifi\'ee si~$X'$ est
quasi-projectif sur~$S'$).
\end{corollaire}

La remarque entre parenth\`eses provient du fait que si~$s$ est le
point de~$S$ au-dessous de~$x'$, alors $R(x')$ est fini et contenu
dans la
fibre~$X'_{s}$,
d'autre part comme~$X'$ est quasi-projectif sur~$S'$ et~$S'$ est fini
sur~$S$, $X'$ est quasi-projectif sur~$S$, ce qui implique qu'une
fibre
de~$X'$ sur~$S$
est contenue dans un ouvert affine.

Comme toute partie finie d'un sch\'ema affine admet un syst\`eme
fondamental de voisinages affines, on voit qu'on ne perd pas
l'hypoth\`ese en se restreignant au-dessus d'un ouvert affine
de~$S$, ce qui en vertu de~\ref{VIII.7.3} nous ram\`ene au cas
o\`u~$S$ est affine. En vertu de~\ref{VIII.7.4}, on est ramen\'e
\`a montrer
que~$x'$
est contenu
dans un ouvert affine \emph{stable} par la donn\'ee de
descente. Soit en effet~$U$ un ouvert affine contenant~$R(x')$, alors
le satur\'e
$$
R(X'-U)=q_{2}(q_{1}^{-1}(X'-U))
$$
ne rencontre pas~$R(x')$, d'autre part comme~$q_{2}$ est fini (car~$g$
donc~$p_{2}$ l'est) donc ferm\'e, le deuxi\`eme membre est une
partie ferm\'ee de~$X'$. Soit~$U'$ son compl\'ementaire dans~$X'$,
c'est donc un ouvert \emph{saturé} et on a
$$
R(x')\subset U'\subset U
$$
avec~$U$ affine, mais~$U'$ pas affine a priori. Comme une partie
finie~$R(x')$ dans un sch\'ema affine~$U$ a un syst\`eme
fondamental de voisinages affines de la forme~$U_{f}$, on voit,
rempla\c cant~$f$ par sa restriction \`a~$U'$, qu'il existe une
section~$f$ de~$\mathcal{O}_{U}$ telle que:
$$
R(x')\subset U_{f}',\quad U_{f}'\textrm{ est affine}.
$$
Soit alors~$U''=q_{1}^{-1}(U')=q_{2}^{-1}(U')$, d\'esignons encore
par~$q_{1},q_{2}$ les morphismes induits~$U''\to U'$, et
consid\'erons
$$
f'=\Norm_{q_{2}}(q_{1}^*(f))\quoi,
$$
o\`u~$\Norm_{q_{2}}$ d\'esigne la \emph{norme} relativement au
morphisme fini localement libre~$q_{2}\colon U''\to U'$. La
compatibilit\'e de la formation de la norme avec le changement de
base implique facilement que~$f'$ est une section \emph{invariante}:
$$
q_{1}^*(f')=q_{2}^*(f')
$$
ce qui implique que~$U_{f'}'$ est un ouvert satur\'e de~$U'$. De
fa\c con plus pr\'ecise d'ailleurs, d\'esignant par~$Z(f')$
l'ensemble des z\'eros d'une section~$f'$, on trouve en vertu des
propri\'et\'es des normes:
$$
Z(f')=q_{2}(Z(q_{1}^*(f)))=q_{2}(q_{1}^*(Z(f)))=
R(U'-U_{f}')\quoi,
$$
ce
qui implique que~$U_{f'}'=U'-Z(f')$ est satur\'e, contient~$R(x')$,
et est contenu dans~$U_{f}'$. Comme ce dernier est affine, il s'ensuit
que~$U_{f'}'$ l'est aussi (car \'egal \`a~$(U_{f}')_{f''}$,
avec~$f''=f'|U_{f'}'$). C'est donc un ouvert affine satur\'e
contenant~$R(x')$ donc~$x'$, ce qui ach\`eve la d\'emonstration.

On notera que ce raisonnement s'applique chaque fois qu'on a une
relation d'\'equivalence (ou m\^eme seulement de
pr\'e\'equivalence, \cf \cite{VIII.3}) dans un pr\'esch\'ema~$X'$, fini
et localement libre et d'ailleurs~\ref{VIII.7.6} est aussi un cas
particulier du r\'esultat analogue pour les pr\'e\'equivalences
finies et localement libres, \cf \loccit M\^eme remarque
pour~\ref{VIII.7.7} ci-dessous.

On peut aussi, une fois obtenue l'existence d'un ouvert quasi-affine
satur\'e~$U'$ contenant~$x'$, faire appel \`a~\ref{VIII.7.9}
et~\ref{VIII.7.2} ce qui \'evite le recours aux normes.

Notons d'ailleurs que sous les conditions de~\ref{VIII.7.6}, si la
donn\'ee de descente sur~$X'$ est effective, $X'$ provenant de~$X$
sur~$S$, alors le morphisme~$X'\to X$ est fini, localement libre et
surjectif, car d\'eduit de~$g$ par le changement de base~$X\to
S$. Il s'ensuit (EGA~II 6.6.4) que si~$X'$ est quasi-projectif
sur~$S'$ donc sur~$S$, alors~$X$ est quasi-projectif sur~$S$, (un
faisceau inversible relativement ample sur~$X$ \'etant obtenu en
prenant la \emph{norme} d'un faisceau inversible sur~$X'$ relativement
ample sur~$S$, ou sur~$S'$, cela revient au m\^eme). On obtient
ainsi:

\begin{corollaire}
\label{VIII.7.7}
Un morphisme~$g\colon S'\to S$ fini localement libre et surjectif est
un morphisme de descente effective pour la cat\'egorie fibr\'ee
des pr\'esch\'emas quasi-projectifs sur d'autres, \ie pour
tout~$X'$ quasi-projectif sur~$S'$, toute donn\'ee de descente
sur~$X'$ relativement \`a~$g$ est effective, et le
$S$-pr\'esch\'ema descendu~$X$ est quasi-projectif sur~$S$.
\end{corollaire}

\begin{proposition}
\label{VIII.7.8}
Soit~$g\colon S'\to S$ un morphisme fid\`element plat et
quasi-compact. Alors~$g$ est un morphisme de descente effective pour
la cat\'egorie fibr\'ee des pr\'esch\'emas~$Z$ quasi-compacts
sur un pr\'esch\'ema~$T$, munis d'un faisceau inversible ample
relativement \`a~$T$. En particulier, pour tout
pr\'esch\'ema~$X'$ sur~$S'$, muni d'une donn\'ee de descente
relativement \`a~$g\colon S'\to S$, et tout faisceau
inversible~$\mathcal{L}'$ sur~$X'$ ample relativement \`a~$S'$, muni
\'egalement d'une donn\'ee de descente relativement \`a celle
donn\'ee sur~$X'$, (\ie muni d'un isomorphisme
de~$q_{1}^*(\mathcal{L}')$ avec~$q_{2}^*(\mathcal{L}')$,
satisfaisant la condition de transitivit\'e habituelle), la
donn\'ee de descente sur~$X'$ est effective, et le faisceau
inversible~$\mathcal{L}$ sur le pr\'esch\'ema descendu~$X$,
d\'eduit de~$\mathcal{L}'$ par descente, est ample relativement
\`a~$S$.
\end{proposition}

La d\'emonstration est toute analogue \`a celle de~\ref{VIII.5.8},
en notant que sur l'Alg\`ebre gradu\'ee
quasi-coh\'erente~$\mathcal{S}'$ sur~$S'$ d\'efinie
par~$\mathcal{L}'$, il y a une donn\'ee de descente, permettant de
construire une Alg\`ebre gradu\'ee
quasi-coh\'erente~$\mathcal{S}$ sur~$S$ gr\^ace
\`a~\ref{VIII.1.1} d'o\`u un~$P=\Proj(\mathcal{S})$ sur~$S$ tel
que~$P'=\Proj(\mathcal{S}')$ s'identifie avec sa donn\'ee de
descente \`a~$P\times_{S}S'$. Comme par hypoth\`ese $X'$
s'identifie \`a un ouvert de~$P'$, n\'ecessairement stable par la
donn\'ee de descente sur~$P'$, la donn\'ee de descente sur~$X'$
est \'egalement effective, et on obtient le pr\'esch\'ema
descendu comme un ouvert dans~$P$. Le d\'etail est laiss\'e au
lecteur. --- En particulier, faisant~$\mathcal{L}'=\mathcal{O}_{X'}$, on
trouve:

\begin{corollaire}
\label{VIII.7.9}
Soit~$g\colon S'\to S$ un morphisme fid\`element plat et
quasi-compact, et soit~$X'$ un pr\'esch\'ema \emph{quasi-affine}
au-dessus de~$S'$, alors toute donn\'ee de descente sur~$X'$
relativement \`a~$g$ est effective, et le pr\'esch\'ema
descendu~$X$ est quasi-affine sur~$S$.
\end{corollaire}

En vertu de~\ref{VIII.6.2}, ce r\'esultat s'applique en particulier
si~$S'$ est localement noeth\'erien et~$X'$ est quasi-fini et
s\'epar\'e sur~$S'$, plus g\'en\'eralement si~$S'$ est
quelconque et~$X'$ est de pr\'esentation finie, quasi-fini et
s\'epar\'e sur~$S'$ (\cf \ref{VIII.6.6}).

\begin{remarques}
\label{VIII.7.10}
Les r\'esultats donn\'es dans ce num\'ero \'epuisent les
crit\`eres actuellement connus d'effectivit\'e, et probablement
m\^eme les crit\`eres utiles existants\footnote{Cette opinion
s'est trouv\'ee partiellement erron\'ee, voir p.\, ex. J.P.\, \textsc{Murre},
S\'em. Bourbaki 294 (Appendix), Mai 1965, et des r\'esultats
sp\'eciaux (notamment de \textsc{Néron} et \textsc{Raynaud}). Pour la descente des
sch\'emas en groupes; \cf M.\, Raynaud, th\`ese (1968).}. On notera
les contre-exemples suivants \`a l'appui de cette assertion:
\begin{enumerate}
\item[(i)]
Si~$S$ est le spectre d'un corps, et~$S'$ est le spectre d'une
extension quadratique galoisienne, on peut trouver un~$X'$ sur~$S'$
propre et lisse sur~$S'$, de dimension 3, muni d'une donn\'ee de
descente qui n'est pas effective (Serre).
\item[(ii)] On peut trouver un~$S$ spectre d'un anneau local
r\'egulier de dimension 3 (si on veut, l'anneau local d'un
sch\'ema alg\'ebrique sur un corps de caract\'eristique
donn\'ee), un~$T$ rev\^etement principal de~$S$ de
groupe~$\ZZ/2\ZZ$, tel que, si~$t$ d\'esigne l'un des points de~$T$
au-dessus du point ferm\'e~$s$ de~$S$, et $S'=T-s$, on puisse
trouver un~$X'$ \emph{projectif} sur~$S'$, r\'egulier, muni d'une
donn\'ee de descente relativement \`a~$g\colon S'\to S$, cette
donn\'ee de descente n'\'etant pas effective.
\end{enumerate}

On utilise pour ces constructions l'exemple de Hironaka de
vari\'et\'es non projectives. Pour~(i), il suffit d'utiliser le
fait qu'on peut trouver au-dessus de~$k$ un sch\'ema propre et
lisse~$X_{0}$ de dimension 3, sur lequel~$G=\ZZ/2\ZZ$ op\`ere sans
inertie, et dans lequel il existe deux points~$a,b$ rationnels
sur~$S$,
congrus sous~$G$, qui ne sont pas contenus dans un ouvert affine. On
pose alors~$X'=X_{0}\times_{k}k'$, on fait op\'erer~$G$ sur~$X'$
gr\^ace aux op\'erations de~$G$ sur les deux facteurs, ce qui
donne une donn\'ee de descente sur~$X'$ relativement \`a~$g\colon
\Spec(k')\to \Spec(k)$. Au-dessus de~$a$ \resp $b$, il y a exactement
un point~$a'$ \resp $b'$, (\`a extension r\'esiduelle
quadratique), et~$a'$ et~$b'$ sont congrus sous~$G$, puisque~$X'\to
X_{0}$ est compatible avec les op\'erations de~$G$. Alors~$a'$
et~$b'$ ne peuvent \^etre contenus dans un ouvert affine, soit~$U'$,
car alors~$U=X_{0}-\Im(X'-U')$ serait un ouvert de~$X_{0}$
contenant~$(a,b)$ et dont l'image inverse dans~$X'$ serait contenue
dans~$U'$, donc quasi-affine, donc~$U$ serait quasi-affine, et par
suite~$(a,b)$ aurait un voisinage affine dans~$U$.

Pour~(ii), on utilise le fait que dans l'exemple de Hironaka, $X_{0}$
est obtenu comme pr\'esch\'ema propre au-dessus
d'un~$k$-sch\'ema projectif~$Y$, lisse sur~$k$ (le
morphisme~$f\colon X_{0}\to Y$ \'etant d'ailleurs birationnel, mais
peu importe), le groupe~$G$ op\'erant \'egalement sur~$Y$ de
fa\c con compatible avec ses op\'erations sur~$X_{0}$, enfin
posant~$S'=Y-f(b)$, $X'=X_{0}|S'$, $X'$ est projectif
sur~$S'$. Alors~$X_{0}$ est muni d'une donn\'ee de descente
naturelle relative au morphisme
canonique~$Y\to S=Y/G$, gr\^ace aux op\'erations de~$G$
sur~$X_{0}$ compatibles avec ses op\'erations sur~$Y$. Cette
donn\'ee de descente n'est pas effective, puisque~$(a,b)$ n'est pas
contenu dans un ouvert affine. La donn\'e de descente induite
sur~$X'$ relativement \`a~$g\colon S'\to S$ n'est alors pas
effective, comme on v\'erifie facilement.
\end{remarques}

\begin{thebibliography}{D}{VIII.8}

\bibitem[D]{VIII.D} J.\, \textsc{Giraud}, M\'ethode de la descente, M\'emoire \no 2
de la Soc.\ Math.\ Fr., 1964.

\bibitem[1]{VIII.1} A.\, \textsc{Grothendieck}, S\'eminaire Bourbaki:
G\'eom\'etrie formelle et G\'eom\'etrie alg\'ebrique, Mai
1959, \No 182.

\bibitem[2]{VIII.2} A.\, \textsc{Grothendieck}, S\'eminaire Bourbaki: Technique de
descente et Th\'eor\`emes d'existence~I, D\'ecembre 1959,
\No 190.

\bibitem[3]{VIII.3} A.\, \textsc{Grothendieck}, S\'eminaire Bourbaki: Technique de
descente et Th\'eor\`emes d'existence~III, F\'evrier 1961,
\No 212.
\end{thebibliography}

